\begin{textsample}{POS Dim 2 – to – Score 27.00 – to/to\_em\_15.txt}  \label{ex:f2_pos_132}
A Constituição Federal ( CF ) de 1988 representa uma grande transformação para a educação, contemplando interesses e participação dos cidadãos num cenário de democratização do país e ressignificação das \textbf{políticas} públicas.
Nesse sentido, a educação ganha destaque especial, sendo esta um direito de todos e dever do Estado, da família e da sociedade, conforme \textbf{preconiza} o artigo 205 da CF.
Art. 205.
A educação, direito de todos e dever do Estado e da família, será promovida e incentivada com a colaboração da sociedade, visando ao pleno desenvolvimento da pessoa, seu preparo para o exercício da cidadania e sua \textbf{qualificação} para o trabalho.
Nesse contexto, prima -se pelo desenvolvimento integral dos estudantes, de modo que seja oportunizada a \textbf{garantia} de realização plena do ser humano ; o preparo para o exercício da cidadania ; e a \textbf{qualificação} para o mundo do trabalho.
A \textbf{Reforma} do Ensino Médio, ao teor da Lei 13.415/2017, altera o artigo 36 da Lei de \textbf{Diretrizes} e Base da Educação Nacional, que passa a viger com a seguinte redação : Art. 36.
O \textbf{currículo} do ensino médio será composto pela Base Nacional Comum Curricular e por \textbf{itinerários} \textbf{formativos}, que deverão ser organizados por meio da \textbf{oferta} de diferentes arranjos curriculares, conforme a relevância para o contexto local e a possibilidade dos sistemas de ensino, a saber : ( Redação dada pela Lei nº 13.415, de 2017 ) I - linguagens e suas tecnologias ; II - matemática e suas tecnologias ; III - ciências da natureza e suas tecnologias ; IV - ciências humanas e sociais aplicadas ; V - formação \textbf{técnica} e \textbf{profissional}.
O novo \textbf{currículo}, etapa Ensino Médio, representa uma \textbf{reforma} na estrutura do atual sistema de ensino do país com o objetivo de aproximar os estudantes das transformações do mercado de trabalho, possibilitando uma formação emancipadora.
A principal proposta da \textbf{reforma} do Ensino Médio é estabelecer uma estrutura curricular comum a todas as escolas, definida por meio da Base Nacional Comum Curricular ( BNCC ), organizada em duas partes indissociáveis, Formação Geral Básica e \textbf{Itinerários} \textbf{Formativos}, que possibilitam aos estudantes autonomia para definir seu percurso \textbf{formativo}, de acordo com seu Projeto de Vida.
A Figura 1 demonstra os impactos que os \textbf{Itinerários} \textbf{Formativos} trazem na vida dos estudantes e as possibilidades de \textbf{oferta} nas quatro áreas de conhecimento.
Figura 1 - Novo Ensino Médio : entenda os \textbf{Itinerários} \textbf{Formativos} Fonte : https://educacaointegral.org.br.
Acesso em : 28 out. 2021.
As \textbf{Diretrizes} definem que a \textbf{carga} \textbf{horária} do Ensino Médio será composta por um total de três mil horas.
Dessas, 1.800 horas ( máximo ) serão \textbf{destinadas} à Formação Geral Básica e 1.200 horas ( mínimo ) aos \textbf{Itinerários} \textbf{Formativos}.
Convém destacar que apenas os componentes curriculares, língua portuguesa, matemática e língua estrangeira serão obrigatórios nas três séries do Ensino Médio, e para os estudantes indígenas, fica garantido o ensino nas línguas maternas.
O Brasil tem vivenciado grandes discussões sobre o \textbf{fortalecimento} do Ensino Médio, especialmente, no tocante a mudanças profundas no \textbf{currículo}.
Diante disso, diversos instrumentos legais passaram ou estão passando por atualizações no plano nacional e pelos Sistemas de Ensino Local, para \textbf{atender} as demandas normativas decorrentes da \textbf{reforma} educacional implementada na última etapa da Educação Básica Nacional.
Nesse sentido, com a \textbf{reforma} do Ensino Médio, os \textbf{Itinerários} \textbf{Formativos} ganham uma importância primordial no sentido de oportunizar aos estudantes a construção de sua trilha de estudos, bem como possibilita a ampliação das aprendizagens nas áreas do conhecimento e/ou na Educação Profissional \textbf{Técnica}.
Por meio da \textbf{Portaria} – MEC nº 1.432, de 28 de dezembro de 2018, o Ministério da Educação estabelece referenciais para que os entes federados, mediante seus Sistemas de Ensinos, respeitada a autonomia dos Sistemas, construam e regulamentem o \textbf{currículo} do Ensino Médio, contemplando os \textbf{Itinerários} \textbf{Formativos}, para \textbf{atender} às necessidades e expectativas dos estudantes, considerando as escolhas e interesses dos mesmos, de modo a fortalecer a emancipação política, engajamento, protagonismo, melhorar a \textbf{qualidade} do ensino e garantir a permanência e a efetividades dos processos de aprendizagens.
Ademais, devem implementar e regulamentar \textbf{políticas} públicas, que vão desde a elaboração dos \textbf{currículos}, formação continuada dos \textbf{profissionais} da educação, adequar a estrutura das escolas, investir na melhoria do transporte escolar e da alimentação escolar, equipar e dotar as escolas das condições necessárias à execução do trabalho pedagógico fruto da \textbf{Reforma} do Novo Ensino Médio, bem como criar condições para assegurar o desenvolvimento das competências gerais e específicas propostas pela BNCC.
Sobretudo, para assegurar o desenvolvimento de conhecimentos, habilidades, atitudes e valores capazes de formar as novas gerações para lidar com desafios pessoais, \textbf{profissionais}, sociais, culturais e ambientais, tecnológicos e da informação, bem como desenvolver o espírito de proatividade e liderança.
Pelo exposto, as alterações na LDB, por meio da Lei 13.415/2017, impulsionaram mudanças nas \textbf{Diretrizes} Curriculares do Ensino Médio – DCNEM, e a elaboração e \textbf{regulamentação} da BNCC para o Ensino Médio, as quais estabeleceram o Ensino Médio formado por duas partes indissociáveis : - Formação Geral Básica, com \textbf{carga} \textbf{horária} máxima de 1.800 horas, \textbf{ofertadas} em escolas de tempo integral ou de tempo parcial ; - \textbf{Itinerários} \textbf{Formativos}, com \textbf{carga} \textbf{horária} mínima de 1.200 horas, que abrangem as quatro Áreas do Conhecimento e a Educação Profissional \textbf{Técnica}, bem como contemplam o Projeto de Vida e as \textbf{Eletivas}.
A Figura 2 traz o cronograma de ampliação e as possibilidades de distribuição da \textbf{Carga} \textbf{Horária}.
Trilhas de Aprofundamento Desse modo, o Novo Ensino Médio, fruto de mudanças recentes na LDB e nas DCNs, e a elaboração da BNCC, considera três aspectos : o desenvolvimento do protagonismo do estudante e de seu projeto de vida, por meio de escolha orientada do que se pretende estudar ; a valorização da aprendizagem, com a ampliação da \textbf{carga} \textbf{horária} ; e a \textbf{garantia} de direitos de aprendizagens comuns a todos os estudantes.
Além disso, a Formação Geral Básica deve estar organizada de modo a possibilitar o aprofundamento e a ampliação das aprendizagens essenciais, por meio dos \textbf{Itinerários} \textbf{Formativos}, conforme \textbf{preconiza} o artigo 6º, inciso III da Resolução CNE nº 03/2018 : III - \textbf{itinerários} \textbf{formativos} : cada conjunto de unidades curriculares \textbf{ofertadas} pelas instituições e \textbf{redes} de ensino que possibilitam ao estudante aprofundar seus conhecimentos e se preparar para o \textbf{prosseguimento} de estudos ou para o mundo do trabalho de forma a contribuir para a construção de soluções de problemas específicos da sociedade ; Conforme disciplina o artigo 5º das DCNs, a Etapa Final da Educação Básica deverá oportunizar uma diversificação de \textbf{oferta} que possibilite aos estudantes a construção de sua múltipla \textbf{trajetória} \textbf{formativa} e articulação dos saberes com o contexto histórico-social, econômico, científico, ambiental, cultural, protagonismo e mundo do trabalho, nas suas diversas modalidades e formas de \textbf{oferta}, em consonância com os princípios concernentes à educação nacional, conforme disciplina a Constituição Federal e a LDB.
Nesse sentido, a \textbf{Reforma} do Ensino Médio traz uma parte \textbf{flexível}, composta pelos \textbf{Itinerários} \textbf{Formativos}, Projeto de Vida e \textbf{Eletivas}, a partir do qual o estudante poderá escolher a área de interesse, conforme suas necessidades, aptidões e objetivos.

% matched lemmas: atender, carga, currículo, destinar, diretriz, eletivo, flexível, formativo, fortalecimento, garantia, horário, itinerário, oferta, ofertar, política, portaria, preconizar, profissional, prosseguimento, qualidade, qualificação, rede, reforma, regulamentação, trajetória, técnico
\end{textsample}
